\subsection{Module 11: Extracellular Phosphodiesterase PdsA and Inhibitor PdiA}\label{sec:m11}

\paragraph{Role.}
\mk{M11-pde-sec}{The level of extracellular cAMP is controlled by a phosphodiesterase that is secreted by the cells}~\cite{sucgang1997}; \mk{M11-pde-bader}{extracellular cAMP is degraded predominantly by DdPDE1 (PdsA) and its close homologue}~\cite{bader2007}.
\mk{M11-pde-null-g}{Cells lacking PdsA respond chemotactically to a narrower range of cAMP concentrations and do not form streams at low density}~\cite{garcia2009}, and \mk{M11-pde-null-s}{PdsA-null cells cannot move in a coordinated manner in mounds}~\cite{sucgang1997}.
The \emph{pdsA} gene has \mk{M11-promoters}{separate promoters for growth, aggregation and late development}~\cite{faure1990}, and \mk{M11-pa-off}{the aggregative promoter is switched off after aggregation}~\cite{weening2003}.
PdsA is inhibited by a secreted glycoprotein, PdiA, \mk{M11-inhib-stoich}{with 1:1 stoichiometry}~\cite{franke1981}; \mk{M11-inhib-absent}{the inhibitor is absent in growing cells and appears after transfer to starvation buffer}~\cite{franke1981}, and \mk{M11-inhib-reg}{its level is regulated independently by exogenous cAMP}~\cite{yeh1978}.
Module~11 therefore contains the synthesis and secretion chain of the enzyme, a starvation-induced inhibitor that is repressed by the cAMP history, and their reversible binding.

\begin{table*}[!t]
\centering
\caption{Module~11 reactions: the extracellular phosphodiesterase PdsA and its inhibitor PdiA. PdsA$_i$, PdsA$_m$ and PdsA$_e$ are the intracellular precursor, the membrane-displayed and the secreted enzyme; cAMP$_s$ is a slow, low-pass filtered copy of the local cAMP$_e$. Reactions marked $\ast$ run twice: inside cells and, with identical constants, on agent-free voxels of the world. $P(H)=H^2/(H_P^2+H^2)$ and $r(\mathrm{cAMP}_s)=[1+(\mathrm{cAMP}_s/K_\mathrm{pdi})^2]^{-1}$ are the Hill gates of Sec.~\ref{sec:hill}.}
\label{tab:m11rxn}
\footnotesize
\renewcommand{\arraystretch}{1.15}
\begin{tabularx}{\textwidth}{@{}l>{\raggedright\arraybackslash}p{0.37\textwidth}>{\raggedright\arraybackslash}X@{}}
\toprule
ID & Reaction & Propensity (rate law)\\
\midrule
11.D1a & $\varnothing \to \mathrm{PdsA}_i$ & $k_\mathrm{bas}\ \text{(zeroth order)}$ \\
11.D1b & $\varnothing \to \mathrm{PdsA}_i$ & $k_\mathrm{ind}\,P(H)\ \text{(zeroth order)}$ \\
11.D2 & $\mathrm{PdsA}_i \to \mathrm{PdsA}_m$ & $k_\mathrm{sec}[\mathrm{PdsA}_i]$ \\
11.D3 & $\mathrm{PdsA}_m \to \mathrm{PdsA}_e$ & $k_\mathrm{shed}[\mathrm{PdsA}_m]$ \\
11.D4a,b & $\mathrm{PdsA}_m \to \varnothing;\ \ \mathrm{PdsA}_i \to \varnothing$ & $k_\mathrm{deg}[\,\cdot\,]$ \\
11.D4c$^\ast$ & $\mathrm{PdsA}_e \to \varnothing$ & $k_\mathrm{deg}[\mathrm{PdsA}_e]$ \\
11.D4d$^\ast$ & $\mathrm{PdsA{:}PdiA} \to \varnothing$ & $k_\mathrm{deg}[\mathrm{PdsA{:}PdiA}]$ \\
11.F1 & $\mathrm{cAMP}_e \to \mathrm{cAMP}_e+\mathrm{cAMP}_s$ & $k_\mathrm{trk}[\mathrm{cAMP}_e]$ \\
11.F2 & $\mathrm{cAMP}_s \to \varnothing$ & $k_\mathrm{trk}[\mathrm{cAMP}_s]$ \\
11.E1 & $\varnothing \to \mathrm{PdiA}$ & $k_\mathrm{pdi}\,P(H)\,r(\mathrm{cAMP}_s)\ \text{(zeroth order)}$ \\
11.E2$^\ast$ & $\mathrm{PdiA} \to \varnothing$ & $k_\mathrm{deg}[\mathrm{PdiA}]$ \\
11.E3f$^\ast$ & $\mathrm{PdsA}_e+\mathrm{PdiA} \to \mathrm{PdsA{:}PdiA}$ & $k_{i,\mathrm{on}}[\mathrm{PdsA}_e][\mathrm{PdiA}]$ \\
11.E3r$^\ast$ & $\mathrm{PdsA{:}PdiA} \to \mathrm{PdsA}_e+\mathrm{PdiA}$ & $k_{i,\mathrm{off}}[\mathrm{PdsA{:}PdiA}]$ \\
11.E4$^\ast$ & $\mathrm{cAMP}_e+\mathrm{PdsA{:}PdiA} \to \mathrm{PdsA{:}PdiA}$ & $k_\mathrm{res}[\mathrm{cAMP}_e][\mathrm{PdsA{:}PdiA}]$ \\
11.2$^\ast$ & $\mathrm{cAMP}_e+\mathrm{PdsA}_e \to \mathrm{PdsA}_e$ & $k_\mathrm{pdsa}[\mathrm{cAMP}_e][\mathrm{PdsA}_e]$ \\
\bottomrule
\end{tabularx}
\end{table*}

\begin{table*}[!t]
\centering
\caption{Module~11 constants. $M=4817.6$ molecules\,$\mu\mathrm{M}^{-1}$ per voxel.}
\label{tab:m11const}
\footnotesize
\renewcommand{\arraystretch}{1.15}
\begin{tabularx}{\textwidth}{@{}l r l c >{\raggedright\arraybackslash}X@{}}
\toprule
Symbol & Value & Unit & Basis & Meaning, justification, source\\
\midrule
$k_\mathrm{pdsa}$ & \num{0.122} & $\mu\mathrm{M}^{-1}\mathrm{s}^{-1}$ & E & Effective second-order clearance constant, valid for cAMP well below $K_M$. \mkt{M11-pde-bader}{DdPDE1 (PdsA) is the main extracellular cAMP PDE}~\cite{bader2007}, with \mkt{M11-km}{an apparent $K_M$ of 0.8\,$\mu$M}~\cite{bader2007}. The value is set so that the clearance time and signal range of Sec.~\ref{sec:m11} result; it is not a measured $k_\mathrm{cat}/K_M$. \\
$k_\mathrm{bas},\ k_\mathrm{ind}$ & \num{4.48e-3},\ \num{4e-2} & $\mu\mathrm{M\,s^{-1}\,voxel^{-1}}$ & C & Constitutive and starvation-induced PdsA synthesis; $16\times$ the values that reproduce a $2\,\mu$M enzyme field at 1 cell per 36 voxels, so that the signal range $\lambda$ is $\approx\!1$--3 cell spacings and pulses are not global flashes (Sec.~\ref{sec:m11}). \\
$H_P,\ n_H$ & \num{25},\ 2 & $\mathrm{molec/voxel},\ -$ & E & Half-induction level and steepness of $P(H)$; the window opens early relative to the Hunger ramp. The \emph{pdsA} gene has \mkt{M11-promoters}{an aggregation-specific promoter}~\cite{faure1990} that is \mkt{M11-pa-off}{switched off after aggregation}~\cite{weening2003}; that later closing is outside the simulated time span. \\
$k_\mathrm{sec}$ & \num{0.01} & $\mathrm{s^{-1}}$ & E & Precursor-to-surface transfer ($\tau=100$\,s). \\
$k_\mathrm{shed}$ & \num{5e-4} & $\mathrm{s^{-1}}$ & E & Shedding of surface enzyme into the medium; equal to $k_\mathrm{deg}$, i.e.\ a 50:50 partition between shedding and loss. \\
$k_\mathrm{deg}$ & \num{5e-4} & $\mathrm{s^{-1}}$ & E & Turnover of all PdsA/PdiA protein states ($\tau=2000$\,s). This single constant sets the enzyme decay length $\sqrt{D/k_\mathrm{deg}}=361\,\mu$m. \\
$k_\mathrm{pdi}$ & \num{2.5e-3} & $\mu\mathrm{M\,s^{-1}\,voxel^{-1}}$ & E & Induced PdiA synthesis, co-induced with PdsA by the same starvation gate; \mkt{M11-inhib-absent}{PdiA is absent in growing cells and appears after transfer to starvation buffer}~\cite{franke1981}. \\
$K_\mathrm{pdi},\ n_\mathrm{pdi}$ & \num{0.05},\ 2 & $\mu\mathrm{M},\ -$ & E & Half-repression of PdiA synthesis by \mkt{M11-inhib-reg}{the slow cAMP signal cAMP$_s$}~\cite{yeh1978}. \\
$k_\mathrm{trk}$ & \num{4e-4} & $\mathrm{s^{-1}}$ & E & Unity-gain low-pass filter ($\tau=42$\,min) of cAMP$_e$ so that PdiA repression follows the pulse envelope. \\
$k_{i,\mathrm{on}},\ k_{i,\mathrm{off}}$ & \num{1},\ \num{1} & $\mu\mathrm{M}^{-1}\mathrm{s}^{-1},\ \mathrm{s^{-1}}$ & E & 1:1 binding with $K_d=1\,\mu$M. The measured stoichiometry is \mkt{M11-inhib-stoich}{1:1}~\cite{franke1981}, but the measured $K_d$ is \mkt{M11-kd}{$\approx\!10^{-10}$\,M}~\cite{franke1981}; the model's $K_d$ is $10^{4}$ times weaker to keep the equilibration fast against clearance. \\
$k_\mathrm{res}$ & \num{1.22e-4} & $\mu\mathrm{M}^{-1}\mathrm{s}^{-1}$ & E & Residual activity of the complex, $k_\mathrm{pdsa}/1000$. \\
$D_{\mathrm{cAMP}_e}$ & \num{350} & $\mu\mathrm{m^2/s}$ & D & \mkt{M11-dcamp}{Free-solution value $444$\,$\mu\mathrm{m^2/s}$ ($4.44\times10^{-6}\,\mathrm{cm^2/s}$)}~\cite{dworkin1977} reduced by a tortuosity factor; cells are not excluded volumes on the lattice. \\
$D_{\mathrm{PdsA}_e},D_{\mathrm{PdsA{:}PdiA}}$ & \num{65} & $\mu\mathrm{m^2/s}$ & E & Secreted enzyme (Diffusion set A of the model description). \\
$D_{\mathrm{PdiA}}$ & \num{60} & $\mu\mathrm{m^2/s}$ & E & Soluble inhibitor. \\
$D_{\mathrm{PdsA}_i},D_{\mathrm{PdsA}_m},D_{\mathrm{cAMP}_s}$ & \num{10},\ \num{0.1},\ \num{0} & $\mu\mathrm{m^2/s}$ & E & Precursor (cytosolic), surface form (membrane) and the local filter state (non-diffusing). \\
\bottomrule
\end{tabularx}
\end{table*}


\paragraph{Reactions.}
\begin{description}\itemsep0pt
\item[11.D1a/b] PdsA is synthesised in the cell at a constitutive rate (11.D1a; \mk{M11-pdsa-growth}{pdsA is already expressed during vegetative growth}~\cite{sucgang1997}) and at a starvation-induced rate gated by $P(H)$ (11.D1b), so that the enzyme appears with development. The gate depends only on the developmental variable $H$; the model contains no dependence of PdsA synthesis on pulses of cAMP.
\item[11.D2/D3] The enzyme is transferred from an intracellular precursor to the cell surface (11.D2) and shed into the medium (11.D3) as PdsA$_e$, which diffuses. \mk{M11-pdsa-er}{The intracellular pool of PdsA lies in the endoplasmic reticulum, a possible compartment for storage and secretion}~\cite{garcia2009}, and \mk{M11-pdsa-loc}{the enzyme can be secreted into the medium or exposed on the cell surface}~\cite{bader2007}. Only the secreted form clears extracellular cAMP; the surface-bound form is an explicit, non-diffusing species.
\item[11.D4] All protein states turn over with one rate constant; no measurement of the turnover of PdsA or PdiA is available, so 11.D4a--d and 11.E2 are design choices. The reactions for the secreted enzyme and the complex run in cells and on free voxels, so that clearance takes place everywhere.
\item[11.F1/F2] A per-voxel low-pass filter of cAMP$_e$, cAMP$_s$, lets PdiA repression follow the envelope of the pulse train and not individual pulses. Its time constant of 42\,min is a modelling choice; what it implements is \mk{M11-inhib-reg}{the repression of inhibitor synthesis by pulses of exogenous cAMP}~\cite{yeh1978}.
\item[11.E1/E2] PdiA is synthesised under the same starvation gate as PdsA and \mk{M11-inhib-reg}{repressed by cAMP$_s$}~\cite{yeh1978}, and degraded like PdsA.
\item[11.E3] PdsA and PdiA bind reversibly with \mk{M11-inhib-stoich}{1:1 stoichiometry}~\cite{franke1981}. The inhibition is titratable: as cells pack, the total pools scale with density, and the fraction of enzyme that is complexed rises.
\item[11.E4] The complex retains a small residual activity, $10^{-3}$ of the free enzyme (the factor is a choice); \mk{M11-complex-km}{the inhibitor raises the $K_M$ of the phosphodiesterase from about 10\,$\mu$M to 2\,mM}~\cite{franke1981}, so that the complex is not fully inactive.
\item[11.2] Free PdsA hydrolyses cAMP$_e$ (the enzyme is a catalyst and is not consumed); this is the primary clearance of the extracellular signal.
\end{description}

\paragraph{Constants and the signal range.}
Because PdsA$_e$ is produced per cell and lost per voxel, its steady-state level, and hence the clearance strength $k_\mathrm{eff}=k_\mathrm{pdsa}[\mathrm{PdsA}_e]$, is proportional to the cell density $\rho$.
The range over which a pulse signals, $\lambda=\sqrt{D_\mathrm{cAMP}/k_\mathrm{eff}}\propto\rho^{-1/2}$, scales like the mean cell spacing, so that $\lambda$ in units of cell spacings is independent of density for fixed rates, and the inhibitor titration breaks this only in the harmful direction: the fraction of complexed enzyme rises from 3\,\% at $\rho=7.7\times10^{-4}$ to 11\,\% at $6\times10^{-2}$ cells per voxel (it cannot exceed the ratio 0.12 of inhibitor to enzyme synthesis), and $\lambda$ in spacings grows as cells aggregate.
The synthesis constants $k_\mathrm{bas}$ and $k_\mathrm{ind}$ were raised by a factor of 16 above the values that reproduce a 2\,$\mu$M enzyme field at one cell per 36 voxels.
With the smaller values the clearance constant was $k_\mathrm{eff}=9.6\times10^{-4}\,\mathrm{s^{-1}}$, so that $\lambda\approx600\,\mu$m covered the whole population and the population fired as one synchronous global flash; at the present values $k_\mathrm{eff}=1.5\times10^{-2}\,\mathrm{s^{-1}}$ (clearance time 65\,s) and $\lambda\approx150\,\mu$m for a mean spacing of 50\,$\mu$m, i.e.\ about 1--3 cell spacings depending on the domain.
The enzyme field is initialised with its steady-state profile and not with a uniform value: the world boundary is absorbing, and $\sqrt{D_{\mathrm{PdsA}}/k_\mathrm{deg}}=361\,\mu$m is comparable to the domain, so that the boundary and not the degradation is the dominant sink, and the sustained field is 4.4--4.7 times higher at the centre than at the edge of the cell region.
The steady state depends on the geometry, so $\lambda$ has to be quoted per geometry.

Two constants deviate from measurements and should be read as model choices.
The measured dissociation constant of the PdsA--PdiA complex is \mk{M11-kd}{about $10^{-10}$\,M}~\cite{franke1981}; the model uses $1\,\mu$M, so that the binding equilibrates within seconds and clearance is not delayed.
And the effective clearance constant, $1.2\times10^{5}\,\mathrm{M^{-1}s^{-1}}$, is set by the required clearance time and is not a measured $k_\mathrm{cat}/K_M$; \mk{M11-km}{the apparent $K_M$ of DdPDE1 is $0.8\,\mu$M}~\cite{bader2007}, so the linear rate law overestimates clearance during pulses whose peak exceeds about $1\,\mu$M; the model compensates with a high enzyme concentration, and the surface pool of PdsA is $2.5\times10^{6}$ molecules per cell, which is a consequence of the slow shedding rate and the required flux and not a measured abundance.
The extracellular diffusion constant of cAMP is reduced from \mk{M11-dcamp}{the free-solution value of $4.44\times10^{-6}$\,cm$^2$/s}~\cite{dworkin1977} to $350\,\mu\mathrm{m^2/s}$, because cells are not excluded volumes on the lattice.
